\honest{Trying to Try}{Eliezer Yudkowsky}{2008}

I open with Yoda: ``No! Try not! Do, or do not. There is no try.'' Years ago I thought this was ``Deep Wisdom'' that is ``actually quite stupid.'' Success is not something you can simply choose, and no plan works with probability 1. Now I think Yoda was wiser than I first realized.

The first elementary technique is to keep talk about a thing apart from talk about the word for it. If I say ``I'm going to flip that switch,'' I mean that I will make and carry out a plan that gives the best chance of a flipped switch. The plan need not be certain. In ordinary speech, ``I'm going to try to flip the switch'' means much the same, except that it admits I might fail.

Much of life's challenge is holding ourselves to a high enough standard, and trying to do a thing is often much easier than doing it. So ``I'm going to `try to flip' the switch'' means that I aim at having a plan that might flip it. For a self-modifying AI planning its own planning, this should end at ``a reflective equilibrium.'' Humans are different. ``It's far easier to convince ourselves that we are `maximizing our probability of succeeding', than it is to convince ourselves that we will succeed.'' Almost any effort will pass for our hardest, if trying our hardest is all we are trying to do. I give no evidence for this. \nb{Research on goal setting supports it. Locke and Latham's 2002 review reports that specific, difficult goals ``consistently led to higher performance than urging people to do their best,'' because do-your-best goals ``have no external referent and thus are defined idiosyncratically.''}

Steven Brust says to ask not what you can do but what needs to be done. When you ask ``What can I do?'', your best is whatever costs you nothing: the money in your pocket, minus lunch. What needs to be done may take three times your life savings. Want to try to make a million dollars? Buy a lottery ticket with your last dollar. You tried your best.

Only when you want, ``above all else,'' to actually flip the switch will you actually put in the effort to maximize the chance of it. I state this without a hedge. If all you want is to have maximized your chances with the resources available, the first plan you think of will do, and whatever you put in will be all that was ``available.''

``Don't try your best. Win, or fail. There is no best.'' \nb{The same review names a limit. On complex tasks that people have not yet learned, a goal set on the outcome can get in the way of learning. In one air-traffic-control simulation, people did better when told to do their best than when given an outcome goal; in a later study, a specific, difficult learning goal beat ``do your best.'' ``Win, or fail'' is an outcome goal.}
